\phantomsection
\color{uniblue}\section*{\centering{\large{ПЕРЕЧЕНЬ СОКРАЩЕНИЙ}\color{white!0}<.>}}
\addcontentsline{toc}{section}{Перечень сокращений}
\color{black}
\begin{longtable}{>{\raggedright\arraybackslash}m{2cm}>{\raggedright\arraybackslash}m{0.5cm}>{\raggedright\arraybackslash}m{20cm}}
\endfirsthead\endhead\endfoot\endlastfoot
DSP & -- & Digital Signal Processing; \\
GOOSE & -- & Generic Object-Oriented Substation Event; \\
IRF & -- & Internal Relay Fault; \\
PPS & -- & Pulse Per Second; \\
SPI & -- & Serial Peripheral Interface; \\
АВ & -- & автоматический выключатель; \\
АРН & -- & автоматика регулирования напряжения; \\
АРНТ & -- & автоматика регулирования напряжения (авто)трансформатора; \\
АСУ~ТП & -- & автоматизированная система управления технологическими процессами; \\
АЦП & -- & аналого-цифровой преобразователь; \\
БИ & -- & блок испытательный; \\
ВН & -- & высшее напряжение; \\
ДО & -- & дочерние общества; \\
ДТм & -- & датчик температуры масла; \\
ДУм & -- & датчик уровня масла; \\
ИО & -- & измерительный орган; \\
ИЧМ & -- & интерфейс человек-машина; \\
КС & -- & контроль синхронизма / контрольная сумма; \\
МУ & -- & местное управление; \\
НКУ & -- & низковольтное комплектное устройство; \\
НН & -- & низшее напряжение; \\
НП & -- & нулевая последовательность; \\
ОВ & -- & оперативный вывод; \\
ОЗУ & -- & оперативное запоминающее устройство; \\
ОП & -- & обратная последовательность; \\
ПДС & -- & преобразователь дискретных сигналов; \\
ПЗУ & -- & постоянное запоминающее устройство; \\
ПМ & -- & привод моторный / приводной механизм; \\
ПО & -- & пусковой орган / программное обеспечение; \\
ПС & -- & предупредительная сигнализация; \\
РАС & -- & регистратор аварийных событий; \\
РЗА & -- & релейная защита и автоматика; \\
РПН & -- & устройство регулирования напряжения под нагрузкой; \\
РЭ & -- & руководство по эксплуатации; \\
СВ & -- & секционный выключатель; \\
ТН & -- & трансформатор напряжения; \\
ТТ & -- & трансформатор тока; \\
ФСУ & -- & функциональная схема устройства; \\
ЦОС & -- & цифровая обработка сигналов; \\
ЦП & -- & центральный процессор; \\
ШЭТ & -- & шкаф электротехнический типовой. \\
\end{longtable}